\clearpage\begingroup\SourceLeaf{667b}
\ArithmeticHead{32}{L}{H. GEPPERT, ÜBER GAUSS’ ARBEITEN ZUR MECHANIK UND POTENTIALTHEORIE.}
\renewcommand{\thefootnote}{\arabic{footnote})}\setcounter{footnote}{58}
\begin{center}\textbf{III. Die Anziehung der Ellipsoide.}\par\medskip
\textls[100]{1. \textsc{Gauss}’ Vorgänger.}\end{center}
\noindent\hspace*{\parindent}Das Anziehungsproblem einer kugelförmigen Masse hatte \textsc{Newton} bereits
1687 in seinen \textit{Principia}\footnote{\textsc{J. Newton}, \textit{Philosophiae naturalis principia mathematica}, 1, 1687, sect. 12. 13.} erledigt. Von ihm rührt auch der Satz her, dass
eine Kugelschale auf einen inneren Punkt gar keine Anziehung ausübt und
auf einen äusseren eine solche, die proportional ist zu der Wirkung ihrer im
Mittelpunkt konzentrierten Masse. (Die Gleichheit dieser Anziehungen wurde erst
später von \textsc{Euler} und \textsc{Bernoulli} bewiesen.) Die Astronomie erforderte weiter
die Bestimmung der Anziehung eines homogenen Ellipsoids, vor allem des
Rotationsellipsoids oder Sphäroids. Auch mit dieser, und zwar mit der
Attraktion des Sphäroids auf die Punkte der verlängerten Achse, hat sich
bereits \textsc{Newton} beschäftigt, und er fand auch den Satz, dass eine von
konzentrischen, ähnlichen und ähnlich gelegenen Sphäroiden begrenzte Schale
auf einen inneren Punkt des Hohlraumes keine Anziehung ausübt. Ein
weiterer bedeutender Beitrag rührt von \textsc{Mac-Laurin}\footnote{\textsc{C. Mac-Laurin}, \textit{De causa physica fluxus et refluxus maris}, 1740. \textit{Treatise of fluxions}, 1, 1742,
cap. 14.} her. Er ermittelte auf
\clearpage\SourceLeaf{668}\ArithmeticHead{33}{R}{III. DIE ANZIEHUNG DER ELLIPSOIDE.}
\HistoricalFoot{X 2 Abh. 7.}{5}
\noindent synthetischem Wege die Anziehung eines Sphäroids auf Punkte im Inneren und
auf der eigenen Oberfläche und zeigte für solche, dass die senkrecht zur Rotationsachse
gerichtete Komponente proportional dem Abstande des Aufpunktes von
der Achse, und ebenso die senkrecht zur Äquatorebene gerichtete Komponente
proportional dem Abstand des Aufpunktes vom Äquator ist. Für äussere,
und zwar auf den Achsen liegende Punkte formulierte \textsc{Mac-Laurin} den heute
nach ihm benannten Satz, dass die Anziehungskräfte konfokaler Sphäroide
auf einen äusseren Punkt sich wie deren Massen verhalten. In der Folge
versuchten sich \textsc{d’Alembert} und \textsc{Lagrange}\footnote{\textsc{Lagrange}, \textit{Sur l’attraction des sphéroïdes elliptiques.} Nouveaux Mém. de l’Acad. des Sciences
de Berlin, 1773, S. 121—148; 1775, S. 273—279; Oeuvres, 3, 1869, S. 619—658. \textit{Remarques générales
sur le mouvement de plusieurs corps qui s’attirent mutuellement en raison inverse des carrés des distances.}
Ibidem 1777, S. 155—174; Oeuvres, 4, 1869, S. 401—418.} am analytischen Beweis der \textsc{Mac-Laurin}schen
Sätze; zwar gelang letzterem der analytische Beweis des \textsc{Mac-Laurin}schen
Theorems für Sphäroide, aber nicht die Lösung des Anziehungsproblems
für das allgemeine Ellipsoid. Die \textsc{Lagrange}sche Arbeit ist dadurch
bedeutend, dass in ihr das allgemeine Attraktionsproblem eines beliebigen
Körpers formuliert wird, wenn die Anziehung eine Funktion \(f(r)\) des Abstandes
\(r\) ist, und dadurch, dass zum erstenmal die Transformation der auftretenden
mehrfachen Integrale auf die zweckmässigeren räumlichen Polarkoordinaten
mit dem Aufpunkt als Pol vorgenommen wird, was auf eine Zerlegung des
Raumes in Elementarkegel hinauskommt, wie sie schon \textsc{Mac-Laurin} bei seinen
synthetischen Betrachtungen verwandt hatte. In seiner zweiten Arbeit machte
\textsc{Lagrange} die wichtige Bemerkung, dass die Komponenten der \textsc{Newton}schen
Anziehung sich als partielle Ableitungen einer Funktion
\[
V=\iiint\frac{dm}{r}
\]
\noindent ausdrücken lassen; hier tritt also zum erstenmal die bei \textsc{Lagrange} noch
namenlose, später so genannte Potentialfunktion auf.

Die vollständige Bestimmung der Anziehung eines Sphäroids auf beliebige
äussere und innere Punkte gab erstmalig \textsc{Legendre}\footnote{\textsc{A. M. Legendre}, \textit{Recherches sur l’attraction des sphéroïdes homogènes.} Mém. présentés par divers savants, vol. 10. Paris 1785, S. 411—434.}, während diejenige für
das dreiachsige Ellipsoid von \textsc{Laplace}\footnote{\textsc{P. S. Laplace}, \textit{Théorie des attractions des sphéroïdes et de la figure des planètes.} Mém. de
l’Acad. Royale de Paris, 1782, S. 113—196; Oeuvres 10, S. 349 ff. Mécanique céleste, 2, livre 3, 1799,
S. 3—22. \textsc{Ostwalds} Klassiker, Band 19, 1890, wo auch die einschlägigen Arbeiten von \textsc{Ivory}, \textsc{Gauss},
\textsc{Chasles} und \textsc{Dirichlet} wiedergegeben sind.} herrührt.
\clearpage\SourceLeaf{669}\ArithmeticHead{34}{L}{H. GEPPERT, ÜBER GAUSS’ ARBEITEN ZUR MECHANIK UND POTENTIALTHEORIE.}
Wir wollen kurz die \textsc{Laplace}sche Arbeit besprechen. Zunächst wird die
\textsc{Lagrange}sche Transformation auf räumliche Polarkoordinaten und die Integration
nach \(r\) ausgeführt, sodass es sich nur noch um die Auswertung von
Doppelintegralen handelt. Nun tritt die in der Natur des Problems begründete
Spaltung in die Fälle eines im Innern oder Äusseren gelegenen Aufpunktes
ein. Hat das Ellipsoid die Gleichung \(x^2+my^2+nz^2=k^2\), so sind bei
einem inneren Punkt die Attraktionskomponenten von \(k\) unabhängig, woraus
der erweiterte \textsc{Newton}sche Satz für ellipsoidische Schalen folgt. Die Komponenten
selbst werden durch elliptische Integrale 2. Gattung ausgedrückt, die
sich im Falle des Sphäroids auf die von \textsc{Legendre} gefundenen elementaren
Funktionen reduzieren und auch noch gelten, wenn der Aufpunkt in die
Ellipsoidfläche selbst fällt. Schwieriger ist die Ermittlung der Anziehung
auf einen äusseren Punkt; dazu berechnet \textsc{Laplace} zunächst das Potential \(V\)
im Aufpunkte \(a,b,c\) und findet für \(V\) als Funktion von \(a,b,c,m,n,k\) eine
komplizierte Differentialgleichung, die in seiner Theorie als deus ex machina
wirkt. Der Ausdruck \(v=\frac{V}{M}\), worin \(M\) die Gesamtmasse ist, wird nun nach
fallenden Potenzen von \(a,b,c\) unter Zusammenfassung der Glieder gleicher
Dimension entwickelt:
\[
v=U^{(0)}+U^{(1)}+U^{(2)}+\ldots.
\]
\noindent Die genannte Differentialgleichung vermittelt eine Rekursionsformel für die \(U^{(i)}\),
aus der sich alle \(U^{(i)}\) und mithin auch \(v\) als von \(k\) unabhängig ergeben, womit
der \textsc{Mac-Laurin}sche Satz auf dreiachsige Ellipsoide und beliebige äussere Aufpunkte
erweitert ist. Nun kann das äussere Problem leicht auf das innere zurückgeführt
werden, wozu nur die Ermittlung eines mit dem ursprünglichen konfokalen
Ellipsoids, das durch den Aufpunkt geht, nötig ist, was wiederum die
Auflösung einer kubischen Gleichung erfordert. \textsc{Laplace} erscheint in dieser
Arbeit als der erste systematische Ausbeuter des Potentialbegriffs, und dies
erklärt, dass man später oft ihm an Stelle von \textsc{Lagrange} die Einführung desselben
zuschrieb. So elegant die \textsc{Laplace}sche Lösung ist, so benutzt sie doch
vielfach feinsinnige Kunstgriffe, die die Systematik beeinträchtigen, wie die
\clearpage\SourceLeaf{670}\ArithmeticHead{35}{R}{III. DIE ANZIEHUNG DER ELLIPSOIDE.}
\HistoricalFoot{}{5\textsuperscript{*}}
\noindent Differentialgleichung für \(V\), die Zerlegung \(V=Mv\) und die Reihenentwicklung
für \(v\), für die überdies der Beweis fehlt, dass sie bis an die Oberfläche
heran konvergiert.

Den Wunsch nach einer anderen Lösungsmethode erfüllt die 1809 erschienene
Arbeit von \textsc{Ivory}\footnote{\textsc{James Ivory}, \textit{On the attractions of homogeneous ellipsoids.} Phil. Trans. of the Royal Soc. of
London, 1809. Part II, S. 345—372. \textsc{Ostwalds} Klassiker Band 19.}. Dieser belässt die Attraktionskomponenten in
rechtwinkligen Koordinaten und führt die eine Integration aus, was einer prismatischen
Raumzerlegung an Stelle der von seinen Vorgängern verwandten
Kegelzerlegung entspricht. Hat das Ellipsoid die Halbachsen \(k,k_1,k_2\),
so verwendet \textsc{Ivory} zur Darstellung seiner Oberfläche zum erstenmal eine
Parameterdarstellung:
\[
x=k\cos\varphi;\quad y=k_1\sin\varphi\cos\psi;\quad z=k_2\sin\varphi\sin\psi.
\]
\noindent Sind \(h,h_1,h_2\) die Halbachsen eines durch den Aufpunkt gehenden konfokalen
Ellipsoids, ist also
\[
P:\quad a=h\cos m,\quad b=h_1\sin m\cos n,\quad c=h_2\sin m\sin n,
\]
\noindent und der Punkt
\[
P':\quad a'=k\cos m,\quad b'=k_1\sin m\cos n,\quad c'=k_2\sin m\sin n
\]
\noindent der mit dem Aufpunkt korrespondierende Punkt des ursprünglichen Ellipsoids,
so findet \textsc{Ivory} den Satz, dass sich die senkrecht zu einem Hauptschnitt genommenen
Komponenten der Anziehung, die das erste bezw. zweite Ellipsoid auf
\(P\) bezw. \(P'\) ausübt, wie die Flächeninhalte der entsprechenden Hauptschnitte
verhalten. Dieser Satz gestattet die Zurückführung eines äusseren Aufpunktes
auf einen inneren, also nicht auf der Oberfläche gelegenen, und leistet mithin
mehr als der \textsc{Mac-Laurin}sche Satz. Um das innere Problem zu lösen, entwickelt
\textsc{Ivory} die erste Anziehungskomponente \(A\) nach steigenden Potenzen
von \(a\) in der Form
\[
A=A_1a+A_3a^3+A_5a^5+\ldots,
\]
\noindent worin sich die \(A_3,A_5,\ldots\) aus \(A_1\) vermittels Rekursionsformeln ableiten lassen,
die aus der Potentialgleichung entspringen. Da \(A_1\) von \(b,c\) unabhängig ist, so
findet sich \(A_3=A_5=\ldots=0\) und
\[
A=aA_1=a\iint\frac{2x\,dy\,dz}{(x^2+y^2+z^2)^{\frac{3}{2}}},
\]
\clearpage\SourceLeaf{671}\ArithmeticHead{36}{L}{H. GEPPERT, ÜBER GAUSS’ ARBEITEN ZUR MECHANIK UND POTENTIALTHEORIE.}
\noindent entsprechendes für die anderen Komponenten. Hierin ist das \textsc{Newton}sche Theorem
enthalten. Die Einführung der Parameterdarstellung des Ellipsoids erlaubt
die Zurückführung der letztgenannten Doppelintegrale auf einfache elliptische
Integrale. Gegen diese Methode von \textsc{Ivory} ist der Einwand zu erheben, dass
der Beweis der Konvergenz der benutzten Reihenentwicklung für \(A\) im ganzen
Innern des Ellipsoids fehlt.
\clearpage\endgroup
